

\section{Introduction}

\label{sec:intro}

Machine-learning inference is increasingly outsourced to remote services, as modern models often require more computation, memory, and specialized hardware than clients can practically provision locally~\cite{zhang2024hardware,sheng2023flexgen}. In this setting, the service provider controls model execution, while the client typically observes only the returned prediction. The client therefore has no direct assurance that the specified model was evaluated correctly on the submitted input~\cite{kang2024scaling,qu2025zkgpt}. Verifying the result by independently re-executing the model would undermine the efficiency gains of outsourcing~\cite{kang2024scaling}, while inspecting the provider's execution may require access to proprietary model parameters, intermediate values, or other information that either party wishes to keep confidential. 


Zero-knowledge machine learning (zkML) addresses this problem by allowing the service provider to return a cryptographic proof alongside the prediction  ~\cite{chen2024zkml,kang2024scaling,boyuan2021zen,chen2025zktorch}. The proof attests that the specified model computation was executed correctly, while the client verifies the result without repeating the inference. Depending on how the computation is encoded, sensitive values such as model parameters and client inputs can remain hidden from the verifier. 

Many zkML systems~\cite{chen2024zkml,boyuan2021zen} build on zero-knowledge succinct non-interactive arguments of knowledge (ZK-SNARKs)~\cite{bitansky2012extractable,groth2016size}. ZK-SNARKs provide compact proofs with low verification cost, allowing clients to efficiently validate inference results. Their non-interactive form also allows proofs to be stored and verified later by the client or an independent auditor, without further interaction with the prover. The resulting asymmetry is attractive for verification, but places most of the computational burden on the prover.
To prove an inference, the model computation is encoded as arithmetic constraints in a proving framework such as Halo2~\cite{halo2}. For each inference, the prover constructs a witness containing the private inputs and intermediate values needed to satisfy these constraints, and then executes the proving protocol over that witness. This process can be orders of magnitude slower than ordinary inference and can consume large amounts of memory. For example, zkGPT~\cite{qu2025zkgpt} reports 147.7 seconds to prove the generation of 64 tokens, compared with 94.55 milliseconds for the corresponding model execution on the same machine, with peak prover memory reaching 240.5~GB. 


Furthermore, scalability is limited not only by the cost of generating an individual proof, but also by how proving work is scheduled across concurrent inference requests. Existing zkML systems primarily improve the efficiency and practicality of individual proofs through zk-friendly model transformations and optimized circuit construction~\cite{boyuan2021zen,chen2024zkml}, broader model support through frameworks such as EZKL~\cite{ezkl}, and increasingly through GPU acceleration~\cite{sin2024zkllm}, parallel proof accumulation~\cite{chen2025zktorch}, or decomposition into independently provable partitions~\cite{sanjaya2026zkcomposer}. These techniques reduce prover overhead and expose greater parallelism within individual proofs, but largely leave open the complementary serving problem of coordinating many independent inference requests that contend for finite prover memory and compute resources.


As zkML systems move beyond one-off proof generation, they are likely to face the same basic operating condition as conventional inference services, where requests arrive over time and compete for a finite pool of compute and memory. Unlike ordinary inference, however, each request may expand into multiple heterogeneous proof jobs whose resource demands change across proof-construction stages, making concurrent execution a central systems problem.  
Our characterization in Section\ref{sec:characterization}, across three model families, highlights these challenges. Model partitioning exposes parallel proof jobs and can reduce end-to-end proving time, but the resulting jobs vary substantially in execution time and peak memory demand. Resource usage also changes over the course of proof construction, with later cryptographic stages capable of imposing much higher memory demands than earlier preparation stages. Repeated executions of the same circuit can further reuse expensive setup state, making placement sensitive to where that state is already available.  
Together, these properties make naive concurrency ineffective. Launching
too many jobs can oversubscribe memory and cause failures, while reserving
each job's peak demand for its entire lifetime can leave substantial
capacity idle. A zkML serving runtime must therefore decide both
\emph{where} proof jobs execute and \emph{when} resource-intensive
stages may proceed, while accounting for contention and reusable setup
artifacts.

\textbf{Our Work.}
To address this gap, we introduce \TITLE{}, a distributed proof-serving system for zkML inference under concurrent workloads. \TITLE{} partitions the model's computational graph into separately provable components, turning each inference request into a collection of smaller proof jobs. Before proof generation, it computes the intermediate tensor values required at partition boundaries, allowing the resulting jobs to execute independently across prover machines. Cryptographic commitments bind the inputs and outputs of adjacent partitions so that their proofs collectively establish a consistent end-to-end inference. This decomposition exposes parallelism within each request while also creating proof jobs that can be scheduled across concurrent requests.

To coordinate proof execution, \TITLE{} uses a hierarchical resource-management architecture. A global coordinator assigns each proof job to a prover node, while a local scheduler on that node decides when the job's execution stages can run. Instead of fixing the number of concurrent workers or reserving each job's peak memory for its full lifetime, the local scheduler uses resource profiles from prior executions of the same circuit together with the node's current capacity to admit work dynamically. This allows stages from different requests to overlap when their
high-demand phases do not coincide, while keeping each proof job on a
single node throughout execution and avoiding mid-execution migration.



\TITLE{} also exploits the repeated structure of ML inference. Compiled circuits, proving keys, and other setup artifacts can be reused across executions of the same circuit. Because these artifacts are expensive to generate and may also be costly to transfer between machines, the global coordinator considers both resource feasibility and whether the required state is already present on a candidate node. Previously unseen circuits can still execute from a cold start, while recurring workloads benefit from cached proving state and resource profiles learned from prior executions.Compiled circuits, proving keys, and other setup artifacts can be reused

We implement \TITLE{} as a proof-serving layer around existing zkML
backends, with EZKL~\cite{ezkl} as the primary backend and Halo2~\cite{halo2} as the
underlying proof system. The prototype includes the global coordinator,
node-local resource managers, isolated prover processes, and reusable
artifact storage, while leaving the underlying proving protocol
unchanged. We additionally instrument EZKL and Halo2 to expose execution
stages and resource usage for characterization and online admission
decisions.

We evaluate \TITLE{} across convolutional, generative, and transformer
workloads, first isolating the effect of decomposition and then measuring
parallel, distributed, and concurrent execution. On a single worker,
partitioned execution reduces end-to-end proving latency by up to
\(5.6\times\) relative to monolithic proving. Distributing the resulting
proof jobs across multiple prover nodes further reduces latency by up to
\(69\times\) and lowers peak per-machine memory by up to \(48\times\).
\textbf{[Add concurrent-serving result here once finalized.]}

Taken together, these results show that decomposition alone does not
resolve the systems challenges of scalable zkML proving. Once an
inference is expressed as a collection of independently executable proof
jobs, the runtime must manage differences in job cost, changing resource
demand, reusable proving state, and contention across concurrent
requests. \TITLE{} addresses these concerns jointly by treating
partitioned proof generation as a distributed serving problem rather than
only as a source of parallelism within an individual inference.

\textbf{Contributions.} This paper makes the following contributions:

\begin{itemize}
   \item We characterize partitioned zkML inference as a proof-serving
workload. Across convolutional, generative, and transformer models,
we show that proof jobs differ substantially in runtime and memory
demand, that resource usage changes over the lifetime of a job, and
that repeated executions exhibit predictable resource behavior while
allowing expensive setup artifacts to be reused.

    \item We design and implement \TITLE{}, a distributed proof-serving
    runtime for concurrent zkML inference. \TITLE{} separates
    cluster-wide job placement from node-local resource admission,
    using historical execution profiles, current machine capacity, and
    reusable proving state to decide where proof jobs should run and when
    resource-intensive stages may proceed.

    \item We evaluate \TITLE{} across multiple model families and
    deployment scales. Partitioned execution reduces end-to-end proving
    latency by up to \(5.6\times\) on a single worker, while distributed
    execution achieves up to \(69\times\) lower latency than monolithic
    proving and reduces peak per-machine memory by up to \(48\times\).
    \textbf{[Add final concurrent-serving result here once available.]}
\end{itemize}





